%; whizzy paragraph -pdf xpdf -latex ./whizzypdfptex.sh
%; whizzy-paragraph "^\\\\begin{frame}\\|\\\\emtext"
% latex beamer presentation.
% platex, latex-beamer でコンパイルすることを想定。 

%     Tokyo Debian Meeting resources
%     Copyright (C) 2012 Junichi Uekawa

%     This program is free software; you can redistribute it and/or modify
%     it under the terms of the GNU General Public License as published by
%     the Free Software Foundation; either version 2 of the License, or
%     (at your option) any later version.

%     This program is distributed in the hope that it will be useful,
%     but WITHOUT ANY WARRANTY; without even the implied warreanty of
%     MERCHANTABILITY or FITNESS FOR A PARTICULAR PURPOSE.  See the
%     GNU General Public License for more details.

%     You should have received a copy of the GNU General Public License
%     along with this program; if not, write to the Free Software
%     Foundation, Inc., 51 Franklin St, Fifth Floor, Boston, MA  02110-1301 USA

\documentclass[cjk,dvipdfmx,12pt]{beamer}
\usetheme{Tokyo}
\usepackage{monthlypresentation}

%  preview (shell-command (concat "evince " (replace-regexp-in-string "tex$" "pdf"(buffer-file-name)) "&")) 
%  presentation (shell-command (concat "xpdf -fullscreen " (replace-regexp-in-string "tex$" "pdf"(buffer-file-name)) "&"))
%  presentation (shell-command (concat "evince " (replace-regexp-in-string "tex$" "pdf"(buffer-file-name)) "&"))

%http://www.naney.org/diki/dk/hyperref.html
%日本語EUC系環境の時
\AtBeginDvi{\special{pdf:tounicode EUC-UCS2}}
%シフトJIS系環境の時
%\AtBeginDvi{\special{pdf:tounicode 90ms-RKSJ-UCS2}}

\newenvironment{commandlinesmall}%
{\VerbatimEnvironment
  \begin{Sbox}\begin{minipage}{1.0\hsize}\begin{fontsize}{8}{8} \begin{BVerbatim}}%
{\end{BVerbatim}\end{fontsize}\end{minipage}\end{Sbox}
  \setlength{\fboxsep}{8pt}
% start on a new paragraph

\vspace{6pt}% skip before
\fcolorbox{dancerdarkblue}{dancerlightblue}{\TheSbox}

\vspace{6pt}% skip after
}
%end of commandlinesmall

\title{東京エリアDebian勉強会}
\subtitle{第124回 2015年3月度}
\author{野島貴英}
\date{2015年3月7日}
\logo{\includegraphics[width=8cm]{image200607/openlogo-light.eps}}

\begin{document}

\begin{frame}
\titlepage{}
\end{frame}

\begin{frame}{設営準備にご協力ください。}
会場設営よろしくおねがいします。
\end{frame}

\begin{frame}{Agenda}
 \begin{minipage}[t]{0.45\hsize}
  \begin{itemize}
   \item 注意事項
	 \begin{itemize}
	  \item 写真はセミナールーム内のみ可です。
          \item 出入りは自由でないので、もし外出したい方は、野島まで一声くださいませ。
	 \end{itemize}
   \item 事前課題発表
  \end{itemize}
 \end{minipage} 
 \begin{minipage}[t]{0.45\hsize}
  \begin{itemize}
   \item 最近あったDebian関連のイベント報告
	 \begin{itemize}
	 \item 第123回 東京エリアDebian勉強会
         \item OSC 2015 Tokyo/spring出展
	 \end{itemize}
   \item Debian Trivia Quiz
   \item Raspberry Pi 2 Model B に Debian Jessie / armhf をインストールする
   \item 今後のイベント
   \item 今日の宴会場所
  \end{itemize}
 \end{minipage}
\end{frame}

\section{事前課題}
\emtext{事前課題}
{\footnotesize
\begin{prework}{ $BLnEg(B }
  \begin{enumerate}
  \item Q.hack time$B$K2?$r$7$^$9$+!)(B\\
    A. DDTSS$B$7$^$C$9(B!\\
    http://ddtp.debian.net/ddtss/index.cgi/ja
  \item ($B%*%W%7%g%s(B)Q.$B2?$K$D$$$FJ9$-$?$$!?;22C<T$HOC$r$7$?$$$G$9$+!)(B\\
    A. $B;22C<T$H5;=Q$NOC$r$7$?$$!*(B
  \end{enumerate}
\end{prework}

\begin{prework}{ nametake }
  \begin{enumerate}
  \item Q.hack time$B$K2?$r$7$^$9$+!)(B\\
    A. Debian$B3+H/4D6-%;%C%H%"%C%W(B
  \item ($B%*%W%7%g%s(B)Q.$B$I$3$G:#2s$NJY6/2q$N3+:E$rCN$j$^$7$?$+!)(B\\
    A. $B$=$NB>!#(B
  \end{enumerate}
\end{prework}

\begin{prework}{ NOKUBI Takatsugu }
  \begin{enumerate}
  \item Q.hack time$B$K2?$r$7$^$9$+!)(B\\
    A. OpenCV$B!"(BKAKASHI$B!"(BBlender
  \item ($B%*%W%7%g%s(B)Q.$B$I$3$G:#2s$NJY6/2q$N3+:E$rCN$j$^$7$?$+!)(B\\
    A. $B$=$NB>(B
  \item ($B%*%W%7%g%s(B)Q.$B2?$K$D$$$FJ9$-$?$$!?;22C<T$HOC$r$7$?$$$G$9$+!)(B\\
    A. key sign
  \end{enumerate}
\end{prework}

\begin{prework}{ alohaug }
  \begin{enumerate}
  \item Q.hack time$B$K2?$r$7$^$9$+!)(B\\
    A. $B%/%j!<%s%k!<%`4D6-$G(BPGP$B80:n@.!"(BGnuk$BIuF~(B \& dictoss$B$5$s80=pL>(B
  \item ($B%*%W%7%g%s(B)Q.$B$I$3$G:#2s$NJY6/2q$N3+:E$rCN$j$^$7$?$+!)(B\\
    A. twitter (@tokyodebian)
  \item ($B%*%W%7%g%s(B)Q.$B2?$K$D$$$FJ9$-$?$$!?;22C<T$HOC$r$7$?$$$G$9$+!)(B\\
    A. PGP$B804IM}$N$"$k$"$k%M%?<}=8!#(B
  \end{enumerate}
\end{prework}

\begin{prework}{ kenhys }
  \begin{enumerate}
  \item Q.hack time$B$K2?$r$7$^$9$+!)(B\\
    A. $BL$Dj(B
  \item ($B%*%W%7%g%s(B)Q.$B$I$3$G:#2s$NJY6/2q$N3+:E$rCN$j$^$7$?$+!)(B\\
    A. Debian JP$B$N%a!<%j%s%0%j%9%H(B
  \end{enumerate}
\end{prework}

\begin{prework}{ Roger Shimizu }
  \begin{enumerate}
  \item Q.hack time$B$K2?$r$7$^$9$+!)(B\\
    A. $BL$Dj(B
  \item ($B%*%W%7%g%s(B)Q.$B$I$3$G:#2s$NJY6/2q$N3+:E$rCN$j$^$7$?$+!)(B\\
    A. $B$=$NB>(B
  \item ($B%*%W%7%g%s(B)Q.$B2?$K$D$$$FJ9$-$?$$!?;22C<T$HOC$r$7$?$$$G$9$+!)(B\\
    A. GPG$B%-!<%5%$%s$,9T$o$l$k$G$7$g$&$+!#(B
  \end{enumerate}
\end{prework}

\begin{prework}{ ryo\_s }
  \begin{enumerate}
  \item Q.hack time$B$K2?$r$7$^$9$+!)(B\\
    A. $B=q@RFI$_(B
  \item ($B%*%W%7%g%s(B)Q.$B$I$3$G:#2s$NJY6/2q$N3+:E$rCN$j$^$7$?$+!)(B\\
    A. $BM'C#$dCN$j9g$$$+$iD>@\(B
  \end{enumerate}
\end{prework}

\begin{prework}{ dictoss }
  \begin{enumerate}
  \item Q.hack time$B$K2?$r$7$^$9$+!)(B\\
    A. rasbbery pi 2$B$K(BDebian$B$r%$%s%9%H!<%k$9$k2<D4$Y!#(B
  \item ($B%*%W%7%g%s(B)Q.$B$I$3$G:#2s$NJY6/2q$N3+:E$rCN$j$^$7$?$+!)(B\\
    A. Debian JP$B$N%a!<%j%s%0%j%9%H(B
  \item ($B%*%W%7%g%s(B)Q.$B2?$K$D$$$FJ9$-$?$$!?;22C<T$HOC$r$7$?$$$G$9$+!)(B\\
    A. CPU$B$N%]!<%F%#%s%0$O2?$rCN$C$F$$$kI,MW$,$"$k$N$+65$($F2<$5$$!#(B
  \end{enumerate}
\end{prework}

\begin{prework}{ myokoym }
  \begin{enumerate}
  \item Q.hack time$B$K2?$r$7$^$9$+!)(B\\
    A. deb$B%Q%1!<%8$N:n@.<j=g$r0l$+$i3X$S$?$$$H;W$$$^$9!#(B
  \item ($B%*%W%7%g%s(B)Q.$B$I$3$G:#2s$NJY6/2q$N3+:E$rCN$j$^$7$?$+!)(B\\
    A. Debian JP$B$N%a!<%j%s%0%j%9%H(B
  \item ($B%*%W%7%g%s(B)Q.$B2?$K$D$$$FJ9$-$?$$!?;22C<T$HOC$r$7$?$$$G$9$+!)(B\\
    A. deb$B%Q%C%1!<%8$N:n@.$d8x3+$K$D$$$F(B
  \end{enumerate}
\end{prework}

\begin{prework}{ yy\_y\_ja\_jp }
  \begin{enumerate}
  \item Q.hack time$B$K2?$r$7$^$9$+!)(B\\
    A. DDTSS 
  \item ($B%*%W%7%g%s(B)Q.$B$I$3$G:#2s$NJY6/2q$N3+:E$rCN$j$^$7$?$+!)(B\\
    A. $B$=$NB>(B
  \item ($B%*%W%7%g%s(B)Q.$B2?$K$D$$$FJ9$-$?$$!?;22C<T$HOC$r$7$?$$$G$9$+!)(B\\
    A. DDTSS$B$N%l%S%e!<$N$*4j$$(B
  \end{enumerate}
\end{prework}



}

\section{イベント報告}
\emtext{イベント報告}

\begin{frame}{第123回東京エリアDebian勉強会}

\begin{itemize}
\item 場所はスクウェア・エニックスさんのセミナルームをお借りしての開催でした。
\item 参加者は9名でした。
\item セミナ内容は杉本さんによる「Debian GNU/kFreeBSDにおけるJail構築を試してみた」でした。
\item LTは、今井さんによる「Gnukと私」でした。
\item 残りの時間でhack timeを行い、成果発表をしました。
\item 「世界のやまちゃん 新宿花園店」で、久々に宴会をやりました。
\end{itemize} 
  
\end{frame}

\begin{frame}{第123回東京エリアDebian勉強会(つづき)}

  今回セミナ発表は、Debian GNU/kFreeBSD上にJail環境でkFreeBSDや、FreeBSDを稼働させるお話でした。また、linux-i386のシステムを、FreeBSDのLinuxバイナリ互換機能で動作させる試みについても発表がありました。

  LTは、発表者の今井さんがGnuk Tokenを使うに至った動機と、Gnuk Tokenの現状に関する発表でした。Gnuk Token現役利用者の発表は珍しいため非常に興味深い発表でした。
  
\end{frame}

\begin{frame}{OSC 2015 Tokyo/spring出展}

 　毎年のことですが、今年も2/28(土)にて、OSC 2015 Tokyo/springが開かれました。

 　場所はいつもの明星大学 日野キャンパスでした。

 　セミナは、岩松さんにより「Debian Updates」が行われました。20名程度の方に来ていただけたとのことでした。
  
\end{frame}

\begin{frame}{OSC 2015 Tokyo/spring出展(つづき)}

  今回の出展は、
  
  \begin{itemize}
  \item 店員が2名が立ってピッタリなサイズのブースとなりました。（いつもより狭い)
  \item 64bit ARMのボードである、Hikey Board(\url{https://www.96boards.org/products/hikey/})がDebian稼働状態で展示されました。
　\item Gnuk Tokenの展示も行いました。
  \end{itemize}
  
\end{frame}

\section{Debian Trivia Quiz}
\emtext{Debian Trivia Quiz}
\begin{frame}{Debian Trivia Quiz}

  Debian の常識、もちろん知ってますよね?
知らないなんて恥ずかしくて、知らないとは言えないあんなことやこんなこと、
みんなで確認してみましょう。

今回の出題範囲は\url{debian-devel-announce@lists.debian.org},
\url{debian-news@lists.debian.org} に投稿された
内容などからです。

\end{frame}

\subsection{問題}

%; whizzy-master ../debianmeetingresume201311.tex
% $B0J>e$N@_Dj$r$7$F$$$k$?$a!"$3$N%U%!%$%k$G(B M-x whizzytex $B$9$k$H!"(Bwhizzytex$B$,MxMQ$G$-$^$9!#(B
%

\santaku
{2015/3/3$B$K(BDebConf15$B$N%"%J%&%s%9$,9T$o$l$^$7$?!#%9%]%s%5!<%I$J;22C$NEPO?4|8B$O2?;~$^$G$G$7$g$&!)(B}
{2015/3/7}
{2015/3/15}
{2015/3/29}
{C}
{$B=IGqBe$H?);v$,(BDebConf$B3+:EB&;}$A$H$J$k%9%]%s%5!<%I$J;22CEPO?$N!:@Z$O(B3/29$B$G$9!*!:@Z$N;~9o$O(BUTC$B$J$N$+!"%I%$%D$N%m!<%+%k%?%$%`$+!":Y$+$$;v$,$o$+$i$J$$$N$G!";22C$r8!F$$5$l$F$$$kJ}$O!:@Z$KBP$7$F==J,$KF|Dx$NM>M5$r$b$C$FEPO?$5$l$k;v$r$*$9$9$a$7$^$9!#(BDebConf15$B$O!"%I%$%D(B $B%O%$%G%k%P!<%0$G(B2015/8/15-22$B$N3+:E$G$9!#(BDebCamp$B$O!"(B8/9-8/14$B$H$J$j$^$9!#(B}

\santaku
{2015/3/4$B$K(BDPL$B$N:#G/$NA*=P$K$D$$$F%"%J%&%s%9$,$"$j$^$7$?!#(BDPL$B$NN)8uJd!:@Z$O2?;~!)(B}
{2015/3/4}
{2015/3/9}
{2015/4/1}
{B}
{$BKhG/91Nc$N(BDPL$BA*5s$G$9!#(BDPL$B$NN)8uJd!:@Z$O(B3/9$B$G!"A*5s$O(B4/1$B$+$i9T$o$l$^$9!#:#G/$OC/$,N)8uJd$9$k$N$G$7$g$&$+!)$^$?!"F1;~$K!"(BDebian JP Project$B$K$D$$$F$b!"(B2015$BG/$N2qD9$NN)8uJd<TJg=8$,9T$o$l$F$$$^$9!#(B}

\santaku
{2015/3/1$B$K$F!"(BAraki$B$5$s$h$j!"(B2007$BG/$+$i2TF/$7$F$$$?(Bcdn.debian.net$B$,$=$NLr3d$r=*$(!"$"$k(BFQDN$B$r;X$9$@$1$K$J$k$H$$$&$3$H$,%"%J%&%s%9$5$l$^$7$?!#$3$N(BFQDN$B$O0J2<$N$I$l!)(B}
{ftp.debian.org}
{http.debian.net}
{sources.debain.net}
{B}
{cdn.debian.net$B$O!"(Bapt$B$K$h$k(BDebian$B%Q%C%1!<%8$N<hF@@h$K$D$$$F!"(BDNS$B%/%(%j$NH/9T$5$l$?>l=j$K4p$E$$$F!"%f!<%6$K:G$b6a$$%Q%C%1!<%8%j%]%8%H%j$N%5!<%P$N(BIP$B%"%I%l%9$r(BDNS$B$N%l%3!<%I$GJV5Q$9$k;EAH$_$G$9!#(BAraki$B$5$s$K$h$C$F@_7W!"9=C[!"1?MQ$5$l$F$*$j$^$7$?!#(B2010$BG/$K(Bapt$B$,(BHTTP $B%j%@%$%l%/%H$r%5%]!<%H$G$-$k$h$&$K$J$C$?$?$a!"(Bcdn.debian.net$B$N5!G=$r(BHTTP$B$N%j%@%$%l%/%H5!G=$G<B8=$9$k(Bhttp.debian.net$B$,1?MQ$5$l$F$-$^$7$?!#:#2s!"?k$K!"(Bcdn.debian.net$B$r0zB`$5$;$k%"%J%&%s%9$,9T$o$l$?$H$$$&>u67$G$9!#$b$7!"(Bsource.list$B$K(Bcdn.debian.net$B$N(BFQDN$B$r;XDj$7$F$$$k>l9g$O!"(Bhttp.debian.net$B$KJQ$($^$7$g$&!#(B}




\section{Raspberry Pi 2 Model B に Debian Jessie / armhf をインストールする}
\emtext{Raspberry Pi 2 Model B に Debian Jessie / armhf をインストールする}

\section{今後のイベント}
\emtext{今後のイベント}
\begin{frame}{今後のイベント}
\begin{itemize}
 \item 関西エリアDebian勉強会
 \item 東京エリアDebian勉強会　
\end{itemize}
\end{frame}

\section{今日の宴会場所}
\emtext{今日の宴会場所}
\begin{frame}{今日の宴会場所}
未定
\end{frame}

\end{document}

;;; Local Variables: ***
;;; outline-regexp: "\\([ 	]*\\\\\\(documentstyle\\|documentclass\\|emtext\\|section\\|begin{frame}\\)\\*?[ 	]*[[{]\\|[]+\\)" ***
;;; End: ***
